\documentclass[11pt]{article}
\usepackage[T1]{fontenc}
\usepackage[utf8]{inputenc}
\usepackage[margin=1.2in]{geometry}
\usepackage{mathpazo}
\usepackage{courier}
\usepackage{microtype}
\usepackage{amsmath,amssymb,amsthm,mathtools}
\usepackage{booktabs,array,multirow}
\usepackage{graphicx}
\usepackage{enumitem}
\usepackage{setspace}
\usepackage{xcolor}
\usepackage{tikz}
\usetikzlibrary{arrows.meta,positioning,calc,fit,shapes.geometric}
\usepackage{pgfplots}
\pgfplotsset{compat=1.17}
\usepackage{listings}
\usepackage[authoryear,round]{natbib}
\usepackage[labelfont=bf,font=small]{caption}
\usepackage[colorlinks=true,linkcolor=blue!40!black,citecolor=blue!40!black,urlcolor=blue!40!black]{hyperref}

\setstretch{1.08}
\setlength{\parindent}{1.4em}
\setlength{\parskip}{0.2em}
\widowpenalty=10000
\clubpenalty=10000

\newtheorem{definition}{Definition}[section]
\newtheorem{lemma}{Lemma}[section]
\newtheorem{proposition}{Proposition}[section]
\theoremstyle{remark}
\newtheorem{remark}{Remark}[section]
\newtheorem{example}{Example}[section]

\lstset{basicstyle=\ttfamily\footnotesize,breaklines=true,columns=fullflexible,
  frame=single,framerule=0.3pt,rulecolor=\color{black!40},xleftmargin=0.5em,xrightmargin=0.5em,
  keepspaces=true,showstringspaces=false,commentstyle=\itshape\color{black!60}}

\newcommand{\id}{\mathrm{id}}
\newcommand{\Sig}{\Sigma}
\newcommand{\GG}{\mathcal{G}}
\newcommand{\sh}{\mathrm{sh}}
\newcommand{\Dec}{\mathrm{Dec}}
\newcommand{\Key}{\mathrm{Key}}

\title{\textbf{Drifting Scripts}\\[0.35em]
\large Visual Substitution, Glyph Interpolation, and the Receptive Acquisition of New Orthographies}
\author{Flyxion\\ \small Independent Researcher}
\date{October 2026}

\begin{document}
\maketitle

\begin{abstract}
\noindent
A writing system pairs a conventional assignment of marks to linguistic units with a perceptual structure that the assignment leaves untouched, and the second of these governs how easily the first can be changed. This essay develops that distinction through a family of constructions: visual substitution ciphers, in which each letter is replaced by one that resembles it; fonts that separate the stored character from the displayed glyph; and books whose script drifts gradually over the course of a single volume so that a reader may acquire the destination forms receptively while reading content in the ordinary way. The treatment is organised by the status of its claims. Counting results fix the size of the design space of reciprocal alphabets, about $5.3\times10^{14}$ out of $4.0\times10^{26}$ permutations of the Latin capitals. A computed resemblance experiment across four typefaces shows that the best global pairing is strongly font dependent, that only six of thirteen pairs survive in all four fonts, and that a pixel-overlap measure disagrees with intuition about several pairs the design literature would call close. Two geometric results concern intermediate forms: any continuous exchange of two letters by linear interpolation passes through a state in which both letters have the same glyph, and this collapse is removed by a detour through a third direction in shape space, of which Arabic letter dotting is a historical analogue. A scheduling analysis shows that the exposure a reader receives is governed by letter frequency, so that the rarest letter of a model book of eighty thousand words is seen about 278 times in total, and that starting rare letters first attains a lower bound on the length of a drift. The essay closes with formal interpretations in category theory, type theory and sheaf theory, each stated with explicit definitions, including a precise sense in which a drifting alphabet defines a compatible family of local decodings that does not glue to a global key. Successful acquisition is treated throughout as a hypothesis, and an empirical programme is specified to test it.
\end{abstract}

\vspace{0.4em}
\noindent\textbf{Keywords:} writing systems, arbitrariness of the sign, substitution ciphers, variable fonts, perceptual learning, receptive learning, sheaf theory, edition provenance.

\tableofcontents
\newpage

\input{sec1}
\input{sec2}
\input{sec3}
\input{sec4}
\input{sec5}
\input{sec6}
\input{sec7}
\input{sec8}
\input{sec9}
\input{sec10}
\input{sec11}

\appendix
\input{appendices}

\input{bib}
\end{document}
